The collection of all symmetry transformations of a given body is called its 
group of symmetry transformations, or, more briefly, its \emph{symmetry 
group}.  Above, these transformations were described as geometrical motions 
of the body.  In quantum-mechanical applications, however, it is more 
convenient to regard them as transformations of the coordinates under which 
the Hamiltonian of the system remains invariant.  Clearly, if a rotation or 
reflection brings the system into coincidence with itself, the corresponding 
coordinate transformation leaves its Schr\"odinger equation unchanged.  We 
shall therefore speak of the group of transformations under which the given 
Schr\"odinger equation is invariant\footnote{This viewpoint permits us to 
include not only the rotation and reflection groups considered here, but also 
other kinds of transformations that leave the Schr\"odinger equation 
unchanged.  Among them are permutations of the coordinates of identical 
particles belonging to the system (a molecule or an atom).  The collection of 
all permutations of identical particles that are possible in a given system 
is called its \emph{permutation group}; such groups have already occurred in 
\S~63.  The general group properties developed below apply to permutation 
groups as well, but we shall not pursue a more detailed study of groups of 
this kind.

One comment is needed concerning the notation used in this chapter.  Symmetry 
transformations are, in substance, operators of the same kind as those used 
throughout the book, and so their letters ought to carry hats.  We omit the 
hats in deference to customary notation and because no ambiguity can result 
here.  For the same reason the identity transformation is denoted by the 
usual symbol $E$, rather than by 1 as the notation of the other chapters would 
suggest.  Finally, in this chapter the inversion operator is written $I$, in 
place of the symbol $P$ used in \S~30 and customary in present-day quantum-
mechanics literature.}.
